\begin{ouverture}[geometrie]{15}{Programmes de construction}{Partie 4 · Espace et géométrie}
{Perpendiculaires · Parallèles · Cercles · Programmes}
\objectif{Tracer perpendiculaires, parallèles et cercles avec précision.}
\objectif{Exécuter un programme de construction pas à pas.}
\objectif{Élaborer un programme de construction dans des cas simples.}
\prerequis{Instruments et vocabulaire du CM1.}
\end{ouverture}
\pagestyle{otMain}\essentiel\fontsize{11}{15}\selectfont
\phantomsection\label{cours-15-1}
\essSection{1}{Les tracés de référence}
\begin{definition}Deux droites perpendiculaires se coupent en angle droit, vérifié à l’équerre ; deux droites parallèles gardent un écart constant et ne se coupent jamais. Le cercle de centre O et de rayon 3 cm réunit tous les points situés à exactement 3 cm de O.\end{definition}
\begin{savoirfaire}[Équerre et règle ensemble]Perpendiculaire passant par un point : un côté de l’angle droit de l’équerre sur la droite, glisser jusqu’au point, tracer. Parallèle : l’équerre appuyée sur la droite, une règle contre l’équerre, glisser, tracer.\end{savoirfaire}
\begin{exemple}Un diamètre vaut deux rayons : un cercle de rayon $3{,}5$ cm a des diamètres de 7 cm. Le compas reporte aussi les longueurs : il transporte un écartement sans le mesurer.\end{exemple}
\begin{attention}Les codes écrivent les propriétés sur la figure : petit carré pour l’angle droit, traits identiques pour les longueurs égales. Une figure sans codes ne garantit rien.\end{attention}
\clearpage
\phantomsection\label{cours-15-2}
\essSection{2}{Écrire un programme de construction}
\begin{definition}Un programme de construction est une suite d’instructions qui permet de reproduire une figure sans la voir. Chaque phrase nomme les points et les objets : « Trace le segment [AB] de 5 cm. », « Trace le cercle de centre A passant par B. »\end{definition}
\begin{savoirfaire}[Élaborer pas à pas]Pour élaborer un programme : faire un croquis codé ; choisir l’ordre des tracés (par quoi faut-il commencer ?) ; nommer les points au fur et à mesure ; relire en imaginant quelqu’un qui ne voit pas la figure.\end{savoirfaire}
\begin{exemple}Programme du carré ABCD de côté 4 cm : tracer [AB] de 4 cm ; tracer la perpendiculaire à (AB) en A et y reporter AD $=4$ cm ; tracer la perpendiculaire à (AB) en B et y reporter BC $=4$ cm ; joindre D à C. Contrôler le dernier angle.\end{exemple}
\begin{attention}L’ordre des étapes compte : on ne peut pas tracer « le cercle de centre O passant par A » avant d’avoir placé O et A. Un bon programme s’exécute sans retour en arrière.\end{attention}

\clearpage\rubrique{Découverte}
\begin{blocLivre}{ex-15-1}\begin{exercice}[1]{Les bons mots}
Sur la figure, nommer avec la bonne notation : le segment joignant R et S ; la droite passant par T et U ; la longueur du segment [RS].
\label{fig-ex-15-1}\figureGeo{geo-1}
\end{exercice}\end{blocLivre}
\begin{blocLivre}{ex-15-2}\begin{exercice}[1]{Tracer une perpendiculaire}
Tracer la perpendiculaire à la droite (d) passant par le point M.
\label{fig-ex-15-2}\figureGeo{geo-2}
\end{exercice}\end{blocLivre}
\begin{blocLivre}{ex-15-3}\begin{exercice}[1]{Tracer une parallèle}
Tracer la parallèle à la droite (d) passant par le point N.
\label{fig-ex-15-3}\figureGeo{geo-3}
\end{exercice}\end{blocLivre}
\begin{blocLivre}{ex-15-4}\begin{exercice}[1]{Rayon et diamètre}
Un cercle a un rayon de $3{,}5$ cm. Que mesure un diamètre ? Si un cercle a un diamètre de 9 cm, que mesure son rayon ?
\label{fig-ex-15-4}\figureGeo{geo-4}
\end{exercice}\end{blocLivre}
\begin{blocLivre}{ex-15-5}\begin{exercice}[1]{Exécuter un programme}
Exécuter : trace un rectangle ABCD avec AB $=5$ cm et BC $=3$ cm ; place le milieu M de [AB] ; trace le cercle de centre M passant par A.
\label{fig-ex-15-5}\figureGeo{geo-5}
\end{exercice}\end{blocLivre}

\rubrique{Entraînement}
\begin{blocLivre}{ex-15-6}\begin{exercice}[2]{Le triangle du compas}
Exécuter : trace [EF] de 5 cm ; trace l’arc de cercle de centre E et de rayon 4 cm ; trace l’arc de centre F et de rayon 4 cm ; nomme G leur point d’intersection ; trace le triangle EFG. Quelle est la nature du triangle ?
\label{fig-ex-15-6}\figureGeo{geo-6}
\end{exercice}\end{blocLivre}
\begin{blocLivre}{ex-15-7}\begin{exercice}[2]{Le programme en désordre}
Remettre en ordre les instructions du programme du cercle circonscrit au carré : (a) trace le cercle de centre O passant par A ; (b) trace le carré ABCD de côté 4 cm ; (c) trace les diagonales [AC] et [BD] ; (d) nomme O leur point d’intersection.
\label{fig-ex-15-7}\figureGeo{geo-7}
\end{exercice}\end{blocLivre}
\begin{blocLivre}{ex-15-8}\begin{exercice}[2]{L’étape impossible}
Un programme dit : « 1. Trace le cercle de centre O et de rayon 3 cm. 2. Place un point P à 5 cm de O. 3. Le point P est sur le cercle. » Quelle phrase est fausse ? Corriger.
\label{fig-ex-15-8}\figureGeo{geo-8}
\end{exercice}\end{blocLivre}
\begin{blocLivre}{ex-15-9}\begin{exercice}[2]{Écrire le programme}
Écrire un programme de construction pour la figure dessinée : un carré RSTU de côté 3 cm et sa diagonale [RT].
\label{fig-ex-15-9}\figureGeo{geo-9}
\end{exercice}\end{blocLivre}
\begin{blocLivre}{ex-15-10}\begin{exercice}[2]{Le cercle bien défini}
Écrire un programme en deux instructions qui construit un cercle passant par un point A donné, de centre un point O donné. Quelle longueur faut-il reporter au compas ?
\label{fig-ex-15-10}\figureGeo{geo-10}
\end{exercice}\end{blocLivre}

\rubrique{Problèmes}
\begin{blocLivre}{ex-15-11}\begin{exercice}[3]{La figure au téléphone}
Léna dicte à Max : « Trace un losange LMNO de 4 cm de côté. » Max répond : « Il me manque une information pour que nos figures soient superposables. » Qui a raison ? Expliquer.
\label{fig-ex-15-11}\figureGeo{geo-11}
\end{exercice}\end{blocLivre}
\begin{blocLivre}{ex-15-12}\begin{exercice}[3]{Les cercles concentriques}
Exécuter : place un point O ; trace les cercles de centre O et de rayons 1 cm, 2 cm et 3 cm. Quelle est la largeur de chaque anneau obtenu ?
\label{fig-ex-15-12}\figureGeo{geo-12}
\end{exercice}\end{blocLivre}
\begin{blocLivre}{ex-15-13}\begin{exercice}[3]{La cible du stand}
Une cible est formée de trois cercles concentriques de rayons 2 cm, 4 cm et 6 cm. Écrire le programme de construction, puis donner le diamètre du grand cercle.
\label{fig-ex-15-13}\figureGeo{geo-13}
\end{exercice}\end{blocLivre}
\begin{blocLivre}{ex-15-14}\begin{exercice}[3]{La maison à construire}
Exécuter : trace un carré ABCD de côté 4 cm, le côté [DC] en haut ; place le milieu M du segment [DC] ; trace la perpendiculaire à (DC) passant par M et place le point T à 2 cm au-dessus de M ; trace les segments [DT] et [CT]. Quelle figure obtient-on ?
\label{fig-ex-15-14}\figureGeo{geo-14}
\end{exercice}\end{blocLivre}
\begin{blocLivre}{ex-15-15}\begin{exercice}[3]{L’ordre qui compte}
Vrai ou faux : « Dans un programme de construction, on peut toujours échanger deux instructions sans changer la figure. » Justifier avec un exemple.
\label{fig-ex-15-15}\figureGeo{geo-15}
\end{exercice}\end{blocLivre}

\clearpage
